\chapter{振子波函数、自旋二分之一与粒子交换}

\section{无处不在的振子}

任何具有稳定平衡态的系统，在偏离该平衡态时都可以发生振荡。它可能有一个频率，也可能有许多频率。例如，小提琴弦有一个基模和更高的谐波。拨动它会激发这些模式的某种组合。

这里更重要的例子是场。电磁场的平衡构型是真空：没有电场，没有磁场，也没有电磁振荡。将辐射送入空腔或波导，场便可以在其允许的模式中振荡，每个模式都有自己的频率。在量子理论中，这些模式描述光子的集合。

同样的想法也适用于晶体。把每个原子都想象成独立的振子是误导性的。原子之间彼此施加力，因此一个原子的位移会影响其邻居，于是波在晶体中传播。这种耦合运动的简正模式是谐振子，它们的量子称为声子。光子和声子是不同的物理对象，但两者都归结为同样的振子代数。

\section{振子代数}

在振子的讨论中，令质量和约化普朗克常数等于一：
\[
m=1,\qquad \hbar=1.
\]
哈密顿量为
\begin{equation}
H=\frac{p^{2}}{2}+\frac{\omega^{2}x^{2}}{2}.
\end{equation}
参数\(\omega\)确定了势的曲率，同时也是振荡频率。

定义升算符和降算符为
\begin{equation}
A_{\pm}=\frac{p\pm i\omega x}{\sqrt{2\omega}}.
\end{equation}
归一化的选取使得
\begin{equation}
[A_{-},A_{+}]=1.
\end{equation}
数算符为
\begin{equation}
N=A_{+}A_{-}.
\end{equation}
这些方程合起来意味着\(A_{+}\)将\(N\)的一个本征值升高一，而\(A_{-}\)将其降低一。由此得到的谱为
\begin{equation}
n=0,1,2,3,\ldots .
\end{equation}

将两个阶梯算符相乘也显示出对经典能量的量子修正：
\begin{align}
N
&=\frac{1}{2\omega}(p+i\omega x)(p-i\omega x)\\
&=\frac{1}{2\omega}\left(p^{2}+\omega^{2}x^{2}
 +i\omega[x,p]\right)\\
&=\frac{1}{2\omega}\left(p^{2}+\omega^{2}x^{2}-\omega\right),
\end{align}
其中\([x,p]=i\)。因此，
\begin{equation}
H=\omega\left(N+\frac12\right).
\end{equation}
如果附加常数对问题无关紧要，方便的做法是使用平移后的哈密顿量
\begin{equation}
H_{\mathrm{shifted}}=H-\frac{\omega}{2}=\omega N.
\end{equation}
项\(\omega/2\)是振子的零点能，即基态能量。对于场模式，类似的项称为真空能。它反映了不确定性原理：没有任何态能使\(x\)和\(p\)同时严格等于零，因此\(H\)中两个正二次项不可能同时具有零期望值。

\begin{classroomqa}
\lecturetimestamp{00:09:46}
\audiencequestion{为什么数算符的谱中包含零？}
\lecturerresponse{可以通过显式构造基态波函数来证明。该构造将展示一个具有\(N=0\)的态。}
\end{classroomqa}

为了开始这一构造，将态矢量替换为位置的波函数。能量本征态变为
\[
\psi_{0}(x),\ \psi_{1}(x),\ \psi_{2}(x),\ldots,
\]
而基本算符的作用为
\begin{equation}
X\longmapsto x,\qquad
P\longmapsto -i\frac{d}{dx}.
\end{equation}
恢复\(\hbar\)后，第二个替换将是
\(P\mapsto-i\hbar\,d/dx\)。

在求解\(\psi_{0}\)之前，对这个问题有一个更快的抽象答案。在可归一化振子本征态的离散阶梯上，\(N=A_{+}A_{-}\)是非负的，因此以单位步长反复降低不可能无限继续下去。在端点---其可归一化波函数将在下面构造---处，必定存在一个态，记为\(\lvert0\rangle\)，使得
\begin{equation}
A_{-}\lvert0\rangle=\mathbf 0.
\end{equation}
否则\(A_{-}\lvert0\rangle\)将是另一个低一步的态。由此立即得到
\begin{equation}
N\lvert0\rangle
=A_{+}A_{-}\lvert0\rangle
=\mathbf 0
=0\,\lvert0\rangle.
\end{equation}
因此最低态具有数本征值零。

\begin{classroomqa}
\lecturetimestamp{00:14:57}
\audiencequestion{如果零是本征值，那么右侧是否应包含零本征向量？}
\lecturerresponse{右矢 \(\lvert0\rangle\) 是一个归一化的非零向量，其标签为零。在 \(N\lvert0\rangle=0\,\lvert0\rangle\) 中，右侧是零向量。基态右矢与零向量是不同的对象。}
\end{classroomqa}

\section{Schr\"odinger方程与物理解}

现在回到位置表象。不含时的Schr\"odinger方程为
\begin{equation}
\left[
-\frac12\frac{d^{2}}{dx^{2}}
+\frac12\omega^{2}x^{2}
\right]\psi(x)=E\psi(x).
\end{equation}
作为二阶微分方程，它对任意的\(E\)值都有形式解。这些形式解大多不是物理的振子态。它们在\(\lvert x\rvert\)很大时指数增长，因此其范数发散。

对于束缚态，波函数必须平方可积：
\begin{equation}
\int_{-\infty}^{\infty}\psi^{*}(x)\psi(x)\,dx<\infty.
\end{equation}
于是可以将其归一化，使积分等于一。只有特殊的能量值才会产生这样的解。无限直线上的位置态和动量态在这方面是例外，但指数增长的振子解则被直接排除在态空间之外。

\section{高斯基态}

这个二阶方程可以求解，但阶梯代数给出了一个更简单的一阶方程：
\begin{equation}
A_{-}\psi_{0}(x)=0.
\end{equation}
当右端为零时，\(A_{-}\)中的整体因子无关紧要，因此这等价于
\begin{equation}
(p-i\omega x)\psi_{0}(x)=0.
\end{equation}
利用\(p=-i\,d/dx\)得到
\begin{equation}
\left(\frac{d}{dx}+\omega x\right)\psi_{0}(x)=0.
\end{equation}

\begin{classroomqa}
\lecturetimestamp{00:22:32}
\audiencequestion{这能用分离变量法求解吗？}
\lecturerresponse{只有一个自变量，因此没有两个变量可供分离。写成\(\psi_{0}'/\psi_{0}=-\omega x\)，方程具有对数导数的形式。这提示我们把波函数写成指数形式。}
\end{classroomqa}

求解它的一个有用方法是写成
\begin{equation}
\psi_{0}(x)=e^{f(x)}.
\end{equation}
于是
\[
\psi_{0}'(x)=f'(x)e^{f(x)},
\]
公共的指数因子被消去。方程化为
\begin{equation}
\frac{df}{dx}+\omega x=0.
\end{equation}

\begin{classroomqa}
\lecturetimestamp{00:26:28}
\audiencequestion{\(f(x)\)中应该有负号吗？}
\lecturerresponse{是的。由于\(df/dx=-\omega x\)，积分给出一个负的二次项。}
\end{classroomqa}

因此
\begin{equation}
f(x)=-\frac12\omega x^{2}+C,
\end{equation}
以及
\begin{equation}
\psi_{0}(x)\propto e^{-\omega x^{2}/2}.
\end{equation}
乘性归一化常数未加指定。要点在于高斯形状：它光滑、集中在原点附近，并在\(\lvert x\rvert\)很大时迅速衰减。它显然是平方可积的。

无需把答案代回二阶方程来求出其能量。代数运算已经给出
\begin{equation}
E_{0}=\frac{\omega}{2}.
\end{equation}

\section{激发态、节点与经典运动}

第一激发态通过提升基态得到：
\begin{equation}
\psi_{1}(x)\propto A_{+}\psi_{0}(x)
\propto(p+i\omega x)e^{-\omega x^{2}/2}.
\end{equation}
基态方程给出
\[
p\psi_{0}=i\omega x\,\psi_{0}.
\]
改变阶梯算符中的符号，把相消变为相加：
\begin{equation}
(p+i\omega x)\psi_{0}
=2i\omega x\,\psi_{0}.
\end{equation}
不计整体常数，
\begin{equation}
\psi_{1}(x)\propto x\,e^{-\omega x^{2}/2}.
\end{equation}

这个波函数是奇函数：
\begin{equation}
\psi_{1}(-x)=-\psi_{1}(x).
\end{equation}
它的概率密度是偶函数，但在原点处有一个节点，
\begin{equation}
\lvert\psi_{1}(0)\rvert^{2}=0.
\end{equation}
因此第一激发态的概率密度在\(x=0\)处有一个节点。

反复作用\(A_{+}\)会给出乘以同一高斯包络的更高次多项式。\footnote{讲座陈述了更高次多项式因子、每一步新增的节点，以及从\(00{:}34{:}25\)到\(00{:}35{:}49\)的交替对称性。高斯乘以多项式的形式得到Leonard Susskind与Art Friedman的\emph{Quantum Mechanics: The Theoretical Minimum}（2014）中关于振子本征函数讨论的佐证。} 不计归一化，
\begin{equation}
\psi_{n}(x)=P_{n}(x)e^{-\omega x^{2}/2}.
\end{equation}
宇称在偶与奇之间交替，每一步增加一个节点。随着\(n\)增大，波函数向离原点更远处扩展，并出现更快的振荡。

这两个特征有直接的物理意义。离原点更远的支撑区域对应更大的势能，因为势按\(x^{2}\)增长。位置空间中的快速振荡标志着更大的动量。高度激发的振子很可能到达离中心更远处；当它穿过中心区域时，它携带更大的动量，其波函数在那里振荡得更快。

\begin{classroomqa}
\lecturetimestamp{00:37:18}
\audiencequestion{当\(n\)变大时，谐振子是否未能趋近经典行为？}
\lecturerresponse{经典运动不是由一个定态能量本征态来代表的。它出现在包含许多能级的波包中。例如，一个位移后的高斯波包会像经典振子那样来回演化，而且随着摆动幅度增大，相对不确定度变得更小。}
\end{classroomqa}

从基态高斯波包出发，将其从原点移开。被移开的波包不再是能量本征态；它是许多能量本征态的叠加。在含时Schr\"odinger方程下，其各分量获得不同的相位。波包在势的一侧与另一侧之间摆动，在中心附近产生快速的相位变化，而在转折点附近又变得平滑。相对于波包宽度而言，偏移越大，运动看起来就越接近经典。

\begin{classroomqa}
\lecturetimestamp{00:39:29}
\audiencequestion{第二激发态波函数是否应该在原点处穿过零？}
\lecturerresponse{不。它的多项式因子不单纯是\(x^{2}\)；它还包含一个常数项，因此波函数不必在\(x=0\)处为零。}
\end{classroomqa}

阶梯论证还假定\(A_{+}\)和\(A_{-}\)在平方可积的振子定义域上互为厄米共轭，并具有伴随关系所需的边界行为。指数增长的形式解不属于该定义域。

\begin{classroomqa}
\lecturetimestamp{00:40:45}
\audiencequestion{是否存在一种一般性的方法，可以用来判断一个哈密顿量何时能够用这些代数方法求解？}
\lecturerresponse{并不存在那种一般性的规则。精确可解的例子是罕见的。谐振子和理想的非相对论氢原子是代数方法恰好奏效的幸运情形。}
\end{classroomqa}

\section{作为角动量的自旋二分之一}

自旋是粒子独立于其质心轨道运动所携带的角动量。静止的粒子没有轨道角动量\(\mathbf R\times\mathbf P\)，但它仍然可以有自旋。人们可以尝试将其想象为一个微小的旋转物体，或者把自旋当作一种内禀的量子属性。决定性的陈述是数学上的：其分量在空间旋转下按一定方式变换，并满足角动量对易关系。

从这一点起恢复\(\hbar\)。角动量分量满足
\begin{equation}
[L_{z},L_{x}]=i\hbar L_{y},
\end{equation}
另外两个关系由循环置换得到。无量纲的泡利矩阵为
\begin{equation}
\begin{aligned}
\sigma_{z}&=
\begin{pmatrix}
1&0\\
0&-1
\end{pmatrix},\\[0.4em]
\sigma_{x}&=
\begin{pmatrix}
0&1\\
1&0
\end{pmatrix},\\[0.4em]
\sigma_{y}&=
\begin{pmatrix}
0&-i\\
i&0
\end{pmatrix}.
\end{aligned}
\end{equation}
在此基下，\(\sigma_{z}\)是对角的。其特征值为\(+1\)和
\(-1\)。

\begin{classroomqa}
\lecturetimestamp{00:47:35}
\audiencequestion{\(\sigma_x\)中是否有负号？}
\lecturerresponse{没有。\(\sigma_x\)的两个非对角元都是正的。}
\end{classroomqa}

在\(00{:}49{:}35\)处的板书计算开始了逐项
乘法。完成这些乘积得到
\begin{align}
\sigma_{z}\sigma_{x}
&=
\begin{pmatrix}
0&1\\
-1&0
\end{pmatrix}
=i\sigma_{y},\\
\sigma_{x}\sigma_{z}
&=
\begin{pmatrix}
0&-1\\
1&0
\end{pmatrix}
=-i\sigma_{y}.
\end{align}
因此
\begin{equation}
[\sigma_{z},\sigma_{x}]=2i\sigma_{y}.
\end{equation}
多出的因子二通过将物理自旋算符定义为下式而被去掉\footnote{板书计算省略了\(\hbar\)。这里将其恢复，以便\(S_i\)具有角动量的量纲。}
\begin{equation}
S_{i}=\frac{\hbar}{2}\sigma_{i}.
\end{equation}
于是
\begin{equation}
[S_{z},S_{x}]=i\hbar S_{y},
\end{equation}
循环置换的情形类似。因此，恢复\(\hbar\)因子后，泡利矩阵的一半构成一个角动量体系。

令\(\hbar=1\)得到如下所示的形式。\footnote{源帧为\texttt{lecture\_04\_frame\_02.png}。}
\begin{figure}[htbp]
\centering
\includegraphics[width=0.9\linewidth]{lecture_04_frame_02.png}
\caption{在\(\hbar=1\)的单位制下，泡利矩阵的一半就是自旋二分之一的角动量算符，并满足\([S_z,S_x]=iS_y\)。}
\lectureframe{00:55:41}
\end{figure}

两个\(S_z\)本征方程为
\begin{equation}
S_z\lvert\pm\rangle
=\pm\frac{\hbar}{2}\lvert\pm\rangle.
\end{equation}
等价地，无量纲的磁量子数为
\begin{equation}
m_{s}=\pm\frac12,
\qquad
S_{z}\lvert m_{s}\rangle
=m_{s}\hbar\,\lvert m_{s}\rangle.
\end{equation}
角动量升算符和降算符以整数步长移动\(m\)，但这并不要求每条阶梯都经过零。整数多重态包含整数\(m\)值，而半整数多重态包含半整数值。对于自旋二分之一，阶梯只有\(m_s=+1/2\)和\(m_s=-1/2\)两个成员。

这两个态构成自旋二分之一双重态。电子、{\proseunicodefont μ}子、中微子和夸克都属于由这类半整数自旋多重态描述的粒子。

\section{轨道角动量与自旋角动量}

现在有两种角动量。轨道角动量描述相对于原点的运动，
\begin{equation}
\mathbf L=\mathbf R\times\mathbf P,
\end{equation}
而\(\mathbf S\)是内禀自旋。它们的矢量和是总角动量，
\begin{equation}
\mathbf J=\mathbf L+\mathbf S.
\end{equation}

\begin{classroomqa}
\lecturetimestamp{00:56:49}
\audiencequestion{如果一个系统具有总角动量\(\mathbf J\)，如何确定它的自旋\(\mathbf S\)？}
\lecturerresponse{自旋量子数和轨道量子数通常是分别给出的。如果\(L=0\)，则全部角动量都是自旋。否则就需要角动量相加的量子规则，这些规则在此暂缓讨论。}
\end{classroomqa}

这些相加规则同样适用于轨道角动量加自旋、两个自旋，或两个轨道角动量。组合后的系统有其自身允许的总角动量取值，但目前还不需要这些规则。

\begin{classroomqa}
\lecturetimestamp{00:58:18}
\audiencequestion{什么实验证据导致了自旋的引入？}
\lecturerresponse{光谱学提供了证据，但元素周期表给出了一个特别清晰的动机：其壳层结构要求电子具有一种额外的二值性质。}
\end{classroomqa}

\section{类氢简并与壳层计数}

对于轨道角动量，\(L^{2}\)本征值方程为
\begin{equation}
L^{2}\lvert\ell,m\rangle
=\hbar^{2}\ell(\ell+1)\lvert\ell,m\rangle,
\end{equation}
而固定的\(\ell\)值有
\begin{equation}
2\ell+1
\end{equation}
个磁态，对应于\(m=-\ell,-\ell+1,\ldots,\ell\)。

理想的非相对论库仑问题具有额外的简并。其束缚态能量为负，并从下方趋近于零，但计数只取决于能级的对齐。第一个\(\ell=1\)能级与下一个\(\ell=0\)能级对齐；第一个\(\ell=2\)能级并入下一个主壳层；这一模式继续下去。\footnote{\texttt{lecture\_04\_frame\_03.png}，\(01{:}02{:}01\)；完整的主壳层对齐从\(01{:}02{:}33\)到\(01{:}04{:}10\)给出。}

\begin{figure}[htbp]
\centering
\begin{tikzpicture}[x=1.55cm,y=1.0cm]
\draw[->] (-0.65,0) -- (-0.65,4.35);
\draw[->] (-0.85,0) -- (4.75,0);
\node at (-0.95,4.08) {\(E\)};
\node at (4.55,-0.28) {\(\ell\)};
\node at (0,-0.30) {\(0\)};
\node at (2,-0.30) {\(1\)};
\node at (4,-0.30) {\(2\)};

\draw (-0.20,0.75) -- (0.20,0.75);

\draw (-0.20,2.00) -- (0.20,2.00);
\foreach \dx in {-0.18,0,0.18}
  \draw (2+\dx-0.06,2.00) -- (2+\dx+0.06,2.00);

\draw (-0.20,3.15) -- (0.20,3.15);
\foreach \dx in {-0.18,0,0.18}
  \draw (2+\dx-0.06,3.15) -- (2+\dx+0.06,3.15);
\foreach \dx in {-0.32,-0.16,0,0.16,0.32}
  \draw (4+\dx-0.05,3.15) -- (4+\dx+0.05,3.15);

\node[anchor=west] at (0.35,0.75) {\(n=1\)};
\node[anchor=west] at (2.35,2.00) {\(n=2\)};
\node[anchor=west] at (4.40,3.15) {\(n=3\)};
\end{tikzpicture}
\caption{前三个理想库仑壳层合并了具有相同主能量的轨道多重态：第一个壳层中有一个态，第二个中有\(1+3\)个，第三个中有\(1+3+5\)个。}
\lectureframe{01:02:01}
\end{figure}

在主壳层\(n\)中，允许的轨道取值为\footnote{这个紧凑的\(\ell\)范围和下面的求和重构了从\(01{:}02{:}33\)到\(01{:}04{:}10\)所述的能级对齐和计数\(1,4,9,16,\ldots\)。能谱图的起始部分见图像\texttt{lecture\_04\_frame\_03.png}，对应的讲课时刻为\(01{:}02{:}01\)。}
\[
\ell=0,1,\ldots,n-1.
\]
因此空间态的数目为
\begin{equation}
\sum_{\ell=0}^{n-1}(2\ell+1)=n^{2},
\end{equation}
给出序列
\begin{equation}
1,\ 4,\ 9,\ 16,\ldots .
\end{equation}
这只是轨道计数。每个空间态允许两个自旋值，因此相应的理想化总数为
\begin{equation}
2,\ 8,\ 18,\ 32,\ldots=2n^{2}.
\end{equation}

\section{泡利不相容原理与电子的二值性质}

从独立电子近似下的氦开始。其原子核具有两倍质子电荷，这会重新标度类氢能级并将电子向内吸引。暂时忽略电子--电子相互作用，两个电子可以占据最低空间轨道，并给出对氦基态的一个合理的一级近似。

为什么不把第三个电子也放在那里？那会降低它的轨道能量，但这并不会发生。这一限制就是泡利不相容原理：没有两个电子可以占据同一个完整的单粒子态。将其称为普通的排斥力会产生误导。它是对允许量子态的限制。

这里存在一个表面上的矛盾。两个氦电子确实共享同一个空间轨道。解决之道在于，一个完整的电子态所包含的不仅仅是轨道波函数。它还包含一个二值自由度。称这些值为上和下。最低轨道中的一个电子可以是自旋向上，另一个是自旋向下；它们的空间态相同，但它们的完整量子态不同。

在这一点上，计数开始向锂前进，然后又退回来。这一回溯揭示了为什么不相容原理和自旋两者都是必要的。想象一种完全没有自旋的不相容规则。只有一个电子能够占据最低空间态。第二个电子将不得不进入下一个主壳层中的四个空间态之一，从而产生若干种替代方案，而观测到的氦基态却是唯一的。那种无自旋版本的不相容原理是错误的。

恢复两个自旋值。两个电子以相反自旋填满最低轨道。第三个电子，如中性锂中那样，必须进入下一个主壳层。该壳层有一个\(\ell=0\)空间态和三个\(\ell=1\)空间态。在计入自旋之前，有四种可能性。

\begin{classroomqa}
\lecturetimestamp{01:14:58}
\audiencequestion{\(s\)态和\(p\)态如何在第一激发壳层中给出八种可能性？}
\lecturerresponse{有一个\(\ell=0\)的空间态和三个\(\ell=1\)的空间态，共四个空间态。每个允许两个自旋值，所以总数为\(2(1+3)=8\)。}
\end{classroomqa}

这就是元素周期表第二行八个条目背后的简单计数。它有意是一种理想化：后面的讨论将表明，电子--电子相互作用很快就会破坏独立的类氢壳层图像。

这种二值性质也可以作为角动量来检验。电子的自旋伴随着磁矩，因此磁场会使其两个自旋投影具有不同的能量。在零场中原本简并的一个能级分裂为两个，正如一个小磁矩平行或反平行于磁场时所预期的那样。这种分裂检验了新增的二值自由度是否表现为角动量，而不是一个任意的二元标记。

\section{光子、电子与壳层图像的极限}

粒子并非都遵循相同的占据规则。许多光子可以占据同一个电磁模式；一列经典光波可以理解为一个在该模式中包含大量光子的态。电子的行为则不同：泡利不相容原理阻止它们堆积到同一个完整的量子态中。这种对比引出了两种粒子统计之间的区分。

\begin{classroomqa}
\lecturetimestamp{01:20:28}
\audiencequestion{对于自由电子流而非束缚在原子中的电子，情况会怎样？}
\lecturerresponse{不相容规则仍然适用：电子不能全部占据同一个量子态。激光可以把许多光子置于同一个模式中，但电子流无法以这种方式制备。}
\end{classroomqa}

这一限制是电子本身的性质，而不是原子束缚的特殊性。

\begin{classroomqa}
\lecturetimestamp{01:21:12}
\audiencequestion{轨道角动量是否特指绕原子核的旋转？}
\lecturerresponse{角动量是相对于所选原点定义的。在原子中，原子核是自然的原点，而自由粒子相对于所选的任意原点都具有角动量 \(\mathbf R\times\mathbf P\)。}
\end{classroomqa}

每一次关于旋转的讨论都要选择一个旋转所绕的点。原子提供了一个自然的中心，因为其原子核是近似中心势的中心。在自由空间中，没有哪个原点是优先的，不过仍然可以相对于任意选定的原点计算角动量。

\begin{classroomqa}
\lecturetimestamp{01:22:34}
\audiencequestion{为什么元素周期表中会出现重复的八元素行，随后又出现重复的十八元素行？}
\lecturerresponse{简单的独立电子壳层图像无法解释完整的模式。它很快就失效了，因为它忽略了电子之间的相互作用。}
\end{classroomqa}

对于一个大原子，外层电子不仅与原子核相互作用，还与相当可观的内层电子云相互作用。把每个电子都当作在未屏蔽的库仑场中独立运动来处理，已经不再合理。

\begin{classroomqa}
\lecturetimestamp{01:23:33}
\audiencequestion{为什么随着更多电子的加入，原子不会显著变大？}
\lecturerresponse{增加电子的同时也向原子核中增加了质子。增加的核电荷把内层电子向内拉，而这些内层电子又屏蔽了原子核。在这个简单图像中，一个外层价电子随后就在大约为一个单位的有效电荷中运动。}
\end{classroomqa}

增加的核电荷使内层电子云收缩。对于一个只有一个外层价电子的原子，原子核与内层电子在远距离处几乎相互抵消，使价电子处于大致相当于一个单位正电荷的场中。因此，原子的大小并不会简单地与电子数目成正比地增长。

\begin{classroomqa}
\lecturetimestamp{01:24:49}
\audiencequestion{电排斥是相邻原子保持分离的原因吗？}
\lecturerresponse{并非仅靠它。当电子云重叠时，泡利不相容原理提供了一种重要的有效排斥，而对于许多原子来说，这是它们抵抗进一步重叠的主要原因。}
\end{classroomqa}

泡利排斥是不相容原理的一种有效结果，而不是一条额外的力定律。

\begin{classroomqa}
\lecturetimestamp{01:25:23}
\audiencequestion{不相容原理是一条独立的公设吗？}
\lecturerresponse{在非相对论量子力学中，它作为一条公设被引入。在相对论量子理论中，自旋--统计关联为它提供了更深刻的基础。}
\end{classroomqa}

\section{经典力学与量子力学中的全同粒子}

在表述那条公设之前，先考虑经典情况下如何计数全同粒子。把两个同种粒子放入一组盒子中。如果这些粒子暗中带有名字---Harry 和 Sally---，那么 Harry 在第一个盒子中而 Sally 在第二个盒子中，看起来就不同于 Sally 在第一个盒子中而 Harry 在第二个盒子中。

经典力学允许这种图像，即使标记微弱到无法探测。人们可以想象在每个粒子上做一个任意微弱的标记，并仍然把它们的轨迹视为彼此不同。在经典统计力学中，人们通常选择不把全同粒子的置换计为新的位形。对于\(N\)个粒子，这一约定产生熟悉的因子\(1/N!\)。

量子力学更为严格。全同粒子不携带原则上可被追踪的隐藏个体标记。两个电子作为电子是可互换的。把一个电子与一个光子交换会产生物理上不同的安排，因为它们是不同种类，但交换两个电子槽位不能揭示哪个电子事实上是哪个。

单粒子态有一个波函数
\[
\psi(x)=\langle x\vert\psi\rangle.
\]
对于两个粒子，态是两个位置自变量的函数，
\begin{equation}
\psi(x_{1},x_{2}).
\end{equation}
下标标记的是描述中的两个槽位，而不是两个空间方向。其绝对值的平方给出联合位置概率密度。

\section{交换算符与统计假设}

交换两个粒子槽位：
\begin{equation}
\psi(x_{1},x_{2})\longmapsto\psi(x_{2},x_{1}).
\end{equation}
按照本讨论中使用的记号，称该交换算符为\(S\)。这里\(S\)表示交换，而不是自旋矢量\(\mathbf S\)。其作用为
\begin{equation}
(S\psi)(x_{1},x_{2})=\psi(x_{2},x_{1}).
\end{equation}
两个变量的一般函数不必是对称的，因此这一操作可以产生不同的波函数。

交换两次恢复原来的次序：
\begin{equation}
S^{2}=I.
\end{equation}
交换是可逆的，并保持内积，因此它由幺正算符表示。如果\(\psi\)是\(S\)的本征矢量，本征值为\(\lambda\)，那么
\[
S^{2}\psi=\lambda^{2}\psi=\psi,
\]
因此
\begin{equation}
\lambda=\pm1.
\end{equation}

这一结论识别出两个交换本征空间，但它并没有说每个任意的两粒子波函数都已经是\(S\)的本征矢量。任何这样的波函数都可以分解为两部分，\footnote{这一分解是对\(S^{2}=I\)所蕴含的两个本征空间的编辑性代数澄清。所记录的论证陈述了两个本征值，然后引入额外的统计原理。}
\begin{equation}
\psi_{\pm}=\frac12(I\pm S)\psi,
\end{equation}
它们满足
\begin{equation}
S\psi_{\pm}=\pm\psi_{\pm}.
\end{equation}
因此
\begin{align}
\psi_{+}(x_{2},x_{1})&=\psi_{+}(x_{1},x_{2}),\\
\psi_{-}(x_{2},x_{1})&=-\psi_{-}(x_{1},x_{2}).
\end{align}

现在来了一个额外的原理，而不是仅由\(S^{2}=I\)得出的推论。对于任何固定种类的全同粒子，允许的态完全位于这些交换扇区之一中：交换一对粒子总是返回相同的波函数，或者总是返回其负值。\(+1\)扇区是对称的，\(-1\)扇区是反对称的。把这两种可能性物理解释为两种粒子统计是下一步，留待后续讨论。
